\documentclass[12pt]{beamer}
\usetheme{moloch}

\usepackage{amsmath,amscd}
\usepackage{fontspec}
\usepackage{unicode-math}
\usepackage[dvipsnames]{xcolor}

\setbeamercolor{background canvas}{bg=black}
\setbeamercolor{normal text}{fg=white}

\setmonofont{DejaVu Sans Mono}

% Set the title bar to black background with white text
\setbeamercolor{frametitle}{bg=black, fg=white}

% Ensure other structural elements (bullets, etc.) are visible
\setbeamercolor{structure}{fg=white}

\setmainfont{Libertinus Serif}
\setsansfont{Libertinus Sans}
\setmonofont{DejaVu Sans Mono}
\setmathfont{STIXTwoMath-Regular.otf}
%\setmathfont{STIX Two Math}


% Define colors for syntax highlighting
\definecolor{codegreen}{rgb}{0,0.6,0}
\definecolor{codegray}{rgb}{0.5,0.5,0.5}
\definecolor{codecoffee}{rgb}{0.75,0.6,0.3}
\definecolor{question}{rgb}{1.0,0.75,0.1}
\definecolor{lightblue}{rgb}{0.55,0.75,0.9}
\definecolor{yellow}{rgb}{1.0, 1.0, 0.0}
\definecolor{codecomment}{rgb}{0.97, 0.74, 0.4}
\definecolor{codepy}{rgb}{0.9, 0.85, 0.85} %{#F8BC65}
\definecolor{funcColor}{rgb}{0.47, 0.6, 0.8}
\definecolor{midgray}{rgb}{0.4, 0.4, 0.3}
\definecolor{indexwhite}{rgb}{0.9, 0.9, 0.9}

\newcommand{\FrameTitle}[2]{\frametitle{#1 \hfill \small\textnormal{\textcolor{midgray}{#2}}}}
\newcommand{\snl}{\vspace*{4pt}\\}
\newcommand{\nl}{\vspace*{1em}\\}
\newcommand{\vsp}{\vspace*{1em}}
\newcommand{\SixPtVspace}{\vspace*{6pt}}

\newcommand{\Quiz}[1]{{\large\textbf{\color{question}{#1}}}}
\newcommand{\QuizQuestion}[1]{{\Large\textbf{\color{question}{#1}}}}

\newcommand{\Matrix}[1]{\textbf{#1}}
\newcommand{\Vector}[1]{\textbf{#1}}
\newcommand{\MMat}[1]{\mathbf{#1}}
\newcommand{\MVec}[1]{\mathbf{#1}}
\newcommand{\Norm}[1]{\lVert {\mathbf{#1}} \rVert}
\newcommand{\NormBeta}[1]{\lVert \beta \mathbf{#1} \rVert}
\newcommand{\NormConst}[1]{\lVert #1 \rVert}
\newcommand{\Trans}{\text{\raisebox{0.5ex}{\scalebox{0.75}{T}}}}
\newcommand{\abs}[1]{\lvert {#1} \rvert}
\newcommand{\Avg}[1]{\mathbf{avg}(\MVec{#1})}
\newcommand{\Std}[1]{\mathbf{std}(\MVec{#1})}

\begin{document}
\title{\center{Applied Linear Algebra \\\vsp Norm and Distance \\}}

\author{\center{Devi Prasad M \\ Manipal School of Information Sciences \\ Manipal \\}}
\date{\center{August-September, 2026}}
\maketitle

\begin{frame}\FrameTitle{Euclidean Norm}{Unit 1}
The \emph{Euclidean norm} of an $n$-vector $\MVec{x}$, denoted $\Norm{x}$,
is the squareroot of the sum of the squares of its elements,

\begin{align*}
    \Norm{x} & = \sqrt{x^2_1 + x^2_2 + \cdots + x^2_n} \\[8pt]
             & = \sqrt{\MVec{x}^{\Trans} \MVec{x}} \\[8pt]
             & = \Norm{x}_2
\end{align*}

This is also called as the \emph{magnitude} of a vector.\snl

\vspace*{2pt}
We will avoid calling it the \emph{length} of a vector.
\end{frame}

\begin{frame}\FrameTitle{Euclidean Norm}{Unit 1}
    \begin{align*}
        \Norm{\MVec{0}} & = 0 \\[4pt]
        \Norm{e_i} & = 1 \\[4pt]
        \NormConst{[c]} & = \abs{c} \\[4pt]
        \Norm{\MVec{1_n}} & = \sqrt{n} \\[4pt]
    \end{align*}

    we use terms such as a \emph{small} norm and \emph{large} norm
    when we want to compare the magnitudes of vectors.
\end{frame}


\begin{frame}\FrameTitle{Properties of Euclidean Norm}{Unit 1}
    \begin{itemize}
        \item Nonnegative homogeneity. $\NormBeta{x}  = \abs{\beta}\Norm{x}$
        \item Nonnegativity. $\Norm{x} \ge 0$
        \item Definiteness. $\Norm{x} = \MVec{0}$ only if $\MVec{x} = \MVec{0}$.
        \item Triangle inequality. $\Norm{x+y} \le \Norm{x} + \Norm{y}$.
    \end{itemize}
\end{frame}

\begin{frame}\FrameTitle{Properties of Euclidean Norm}{Unit 1}
    \begin{itemize}
        \item Nonnegative homogeneity. $\NormBeta{x}  = \abs{\beta}\Norm{x}$
        \item Nonnegativity. $\Norm{x} \ge 0$
        \item Definiteness. $\Norm{x} = \MVec{0}$ only if $\MVec{x} = \MVec{0}$.
        \item Triangle inequality. $\Norm{x+y} \le \Norm{x} + \Norm{y}$.
    \end{itemize}

    \vspace*{2em}
    \hspace*{1.2em}\Quiz{Quiz:}\snl
    \vspace*{2pt}
    \hspace*{1.2em}{\QuizQuestion{Example for triangle inequality?}}
\end{frame}

\begin{frame}\FrameTitle{Norm - Wikipedia}{Unit 1}
    \emph{a norm is a function from a real or complex vector space
    to the non-negative real numbers that behaves in certain ways like
    the distance from the origin: it commutes with scaling, obeys a form
    of the triangle inequality, and is zero only at the origin.}
\end{frame}


\begin{frame}\FrameTitle{1-\emph{norm} and $\infty$-\emph{norm}}{Unit 1}
    \begin{definition}{General Norm}
        A \emph{general norm} is any real-valued function of an $n$-vector
        that satisfies the four properties listed above.
    \end{definition}

    \vspace*{1em}

    Two examples of general norms:
    \[
    \begin{array}{rll}
        \textbf{1-norm}        &  \lVert x \rVert_1      &= \lvert x_1 \rvert + \cdots + \lvert x_n \rvert \\[4pt]
        \textbf{Infinity-norm} &  \lVert x \rVert_{\infty} &= \max \{\lvert x_1 \rvert, \ldots, \lvert x_n \rvert \
    \end{array}
    \]

    \vspace*{2em}
    Infinity norm $\Norm{\cdot}_{\infty}$ is exrensively used in
    ML-DSA, NIST's PQC signature scheme.
\end{frame}

\begin{frame}\FrameTitle{Root Mean Square Value}{Unit 1}
The norm is related to the root-mean-square (RMS) value of an $n$-vector x:
\begin{align*}
    \mathbf{rms}(\MVec{x}) &= \sqrt{\frac{x^2_1 + \cdots + x^2_n}{n}} \\[4pt]
                    &= \frac{\Norm{x}}{\sqrt{n}}
\end{align*}

The RMS value of a vector $\MVec{x}$ is useful when comparing norms of vectors
with different dimensions.

\vspace*{4pt}
If all the entries of a vector are close to $\alpha$, then the RMS
value of the vector is $\abs{\alpha}$.
\end{frame}

\begin{frame}\FrameTitle{Root Mean Square Value}{Unit 1}
{\color{blue!40}The RMS value tells us what a 'typical' value of $\abs{x_i}$ is.}

\vsp
The norm of $\MVec{1}_n$, the $n$-vector of all ones, is $\sqrt{n}$, but its RMS value is 1,
\emph{independent} of $n$.

\vsp
A couple of \emph{norm identities}
\[
    \begin{array}{rll}
    \Norm{x + y} &=& \sqrt{\Norm{x}^2 + 2 \MVec{x}^{\Trans} \MVec{y} + \Norm{y}^2}\vspace*{6pt}\\
    (\MVec{x} + \MVec{y})^\Trans (\MVec{x} - \MVec{y}) &=& \Norm{x}^2 - \Norm{y}^2
    \end{array}
\]

\end{frame}


\begin{frame}\FrameTitle{Root Mean Square Value}{Unit 1}
    \begin{itemize}
        \item[--] 1805 — Adrien-Marie Legendre. His criterion was
        to make the sum of the squares of the errors as small as possible.
        \item[--] Karl Pearson is generally credited with introducing the
        term “root-mean-square” in 1895, although the underlying
        mathematical quantity, and even root-mean-square deviation,
        predates the terminology.
        \item[--] Pearson's important contribution was to name and
        formalize terminology around it.
        \item[--] In his 1893 lectures, he introduced
        \emph{standard deviation} as a replacement for the older
        \emph{root mean square} error.
    \end{itemize}
\end{frame}

\begin{frame}\FrameTitle{Euclidean Distance}{Unit 1}
    If $\MVec{a}$ and $\MVec{b}$ are $n$-vectors,

    \[
        {\mathbf{dist}}(\MVec{a}, \MVec{b}) = \Norm{\MVec{a} - \MVec{b}}
    \]

    \vspace*{1em}
    The RMS value of the vector difference,
    \emph{RMS deviation} between vectors $\MVec{a}$ and $\MVec{b}$ is
    \[
        \begin{array}{rll}
            \textrm{RMSD} & = & \sqrt{\frac{1}{n}\sum\limits_{i=1}^{n}(a_{i}-b_{i})^{2}} \vspace*{12pt}\\
                 & = & \dfrac{\mathbf{dist}(\MVec{a}, \MVec{b})}{\sqrt{n}}
        \end{array}
    \]
\end{frame}

\begin{frame}\FrameTitle{Euclidean Distance - Example}{Unit 1}
    Given three vectors
    \begin{align*}
    \MVec{u} &= [1.5, -2.5, 1.5, -1.5] \\
    \MVec{v} &= [0.5, -1.5, 3.5, -0.5] \\
    \MVec{w} &= [-1.5, 1.5, 2.5, -2.5]
    \end{align*}

    And the formula
    \[
        \mathbf{dist}(\MVec{x}, \MVec{y}) = \sqrt{(x_1 - y_1)^2 + (x_2 - y_2)^2 + (x_3 - y_3)^2 + (x_4 - y_4)^2}
    \]

    \begin{align*}
        \MVec{u} - \MVec{v} &= [1.0, -1.0, -2.0, -1.0] \\
        \mathbf{dist}(\MVec{u}, \MVec{v}) &= \sqrt{(1.0)^2 + (-1.0)^2 + (-2.0)^2 + (-1.0)^2} \\
                                          &= \sqrt{1.0 + 1.0 + 4.0 + 1.0} \\
                                          &= \sqrt{7} \approx 2.65
    \end{align*}
\end{frame}

\begin{frame}\FrameTitle{Euclidean Distance - Example}{Unit 1}
    \begin{align*}
        \MVec{v} - \MVec{w} &= [2.0, -3.0, 1.0, 2.0] \\
        \mathbf{dist}(\MVec{v}, \MVec{w}) &= \sqrt{(2.0)^2 + (-3.0)^2 + (1.0)^2 + (2.0)^2} \\
                            &= \sqrt{4.0 + 9.0 + 1.0 + 4.0} \\
                            &= \sqrt{18} \approx 4.24
    \end{align*}

    \begin{align*}
        \MVec{u} - \MVec{w} &= [3.0, -4.0, -1.0, 1.0] \\
        \mathbf{dist}(\MVec{u}, \MVec{w}) &= \sqrt{(3.0)^2 + (-4.0)^2 + (-1.0)^2 + (1.0)^2} \\
                            &= \sqrt{9.0 + 16.0 + 1.0 + 1.0} \\
                            &= \sqrt{27} \approx 5.20
    \end{align*}
\end{frame}

\begin{frame}\FrameTitle{Euclidean Distance - An Application(KNN)}{Unit 1}
    \textbf{For Regression (Predicting a Number) - Find the Mean}.\vsp
    \begin{itemize}
        \item \textbf{Measure:} Calculate the Euclidean distance
            between the new query point and \textit{every} point
            in the training dataset.
        \item \textbf{Select:} Identify the $K$ training points
            with the smallest distances (the nearest neighbors).
        \item \textbf{Predict:} Average the \textit{target values}
            (not the distances!) of these $K$ neighbors to
            approximate the query's unknown value.
    \end{itemize}

    \vsp
    Note: Measure and Select are both exhuastive operations.
    \vfill
\end{frame}

\begin{frame}\FrameTitle{Euclidean Distance - KNN}{Unit 1}
    Euclidean distance is the linear-algebraic notion of similarity.
    KNN turns this idea of geometric proximity into an estimate of
    an unknown function assuming that nearby points tend to have
    similar target values.

    \vspace*{1em}
    Linear algebra supplies the geometry, and machine learning supplies
    the statistical assumption that \emph{nearby points or neighbors
    behave similarly}.
    \vspace*{1em}\vfill
\end{frame}

\begin{frame}\FrameTitle{Euclidean Distance - KNN}{Unit 1}
    \noindent\textbf{1. The Distance Between Query and Data Point}\vsp\\

    Let $\MVec{x}_q$ be a new query data point, and $\MVec{x}_i$ be a
    training data point in an $n$-dimensional feature space.

    \SixPtVspace
    The Euclidean distance between these two points is:
    \vspace*{3pt}
    \begin{equation*}
        \mathbf{dist}(\MVec{x}_q, \MVec{x}_i) = \sqrt{\sum_{j=1}^{n} (x_q^{(j)} - x_i^{(j)})^2}
    \end{equation*}

    \vsp
    Let $\mathbf{D} = \{(\MVec{x}_1, y_1), (\MVec{x}_2, y_2), \dots, (\MVec{x}_M, y_M)\}$
    be the training dataset, where $y_i$ is the \emph{known scalar target value}.

    \vsp
    For a given query point $\MVec{x}_q$, let $N_K(\MVec{x}_q) \subset \mathbf{D}$
    represent the subset containing its $K$-closest neighbors.

    \vfill
\end{frame}

\begin{frame}\FrameTitle{Euclidean Distance - KNN}{Unit 1}
    \noindent\textbf{2. The K-NN Regression Prediction Formula} \vsp\\
    Recall $\mathbf{D}$ is the training dataset and $N_K(\MVec{x}_q)$ is the set of $K$-nearest neighbors.

    \vsp
    The predicted value $\hat{y}(\MVec{x}_q)$ is computed as the
    arithmetic mean of the target values of those neighbors:
    \begin{equation*}
        \hat{y}(\MVec{x}_q) = \frac{1}{K} \ \sum_{(\MVec{x}_i, y_i) \in N_K(\MVec{x}_q)} \hspace*{-1.6em}y_i
    \end{equation*}

    \vsp
    \textit{Note: $\hat{y}$ is an approximation of the true unknown value, subject to the neighborhood definition:}
    \begin{equation*}
        N_K(\MVec{x}_q) = \arg\min_{\substack{\vspace*{2pt}\\S \subset \mathbf{D} \vspace*{2pt}\\ |S| = K}} \ \, \sum_{(\MVec{x}_i, y_i) \in S} \hspace*{-0.7em} \mathbf{dist}(\MVec{x}_q, \MVec{x}_i)
    \end{equation*}
    \vfill
\end{frame}

\begin{frame}\FrameTitle{Euclidean Distance - KNN}{Unit 1}
    $N_K(\MVec{x}_q)$ is the set of $K$ training samples nearest to $\MVec{x_q}$.\\

    \[
        N_K(\MVec{x}_q) = \{ \MVec{x}_1, \MVec{x}_2, \ldots, \MVec{x}_K \}
    \]
    where $\mathbf{dist}(\MVec{x}_q, \MVec{x}_1)
        \le \cdots
        \le \mathbf{dist}(\MVec{x}_q, \MVec{x}_K)
        \le \mathbf{dist}(\MVec{x}_q, \MVec{x}_{K+1})$

    \vspace*{2em}
    Euclidean distance is the linear-algebraic notion of similarity.
    KNN turns this idea of geometric proximity into an estimate of
    an unknown function assuming that nearby points tend to have
    similar target values.

    \SixPtVspace
    Linear algebra supplies the geometry, and machine learning supplies
    the statistical assumption that \emph{nearby points or neighbors
    behave similarly}.
\end{frame}


\begin{frame}\FrameTitle{Euclidean Distance - Deviation}{Unit 1}
    \begin{itemize}
        \item When the entries of a vector represent different types of quantities,
            we should carefully choose the units used to represent the numerical
            values of the entries.
        \item  Different components/entries should have approximately equal
                status in determining the vector distance.
        \item The numerical values should be approximately of the same magnitude.
    \end{itemize}
\end{frame}

\end{document}